\section{Rooted costs on invariant measured graphs}\label{sec:invariant}

The preceding finite polynomial estimates control the effect of assigning a
single vertex.  We now pass from finite sets to measured graphs, where many
vertices must be changed simultaneously.  There are two points to establish.
First, the integral of the rooted positive cost changes by the integral of the
single-vertex deletion increments, evaluated in an ordered sequence of
comparison environments.  Second, deleting vertices outside a fixed weight
band has a small effect on the transforms within that band, with a bound
independent of the largest weight anywhere in the graph.  The latter
uniformity will allow us to remove weight cutoffs in the next section.

Let $(X,\nu)$ be a standard finite measure space, and let $\mathcal R$ be a
countable measured equivalence relation preserving $\nu$.  Thus each
measurable partial isomorphism whose graph lies in $\mathcal R$ preserves
measure.  Let $G\subset\mathcal R$ be a measurable undirected graph without
loops, of degree at most $D_0<\infty$.  On its orbit spaces suppose that we are
given measurable Hermitian matrices $T$ of norm at most one, with zero
diagonal and entries supported on $G$.  As before, changing the root may
conjugate these matrices by diagonal unitaries.  All the constructions below
are invariant under those conjugacies.  Null exceptions are discarded along
their $\mathcal R$-saturations.

Fix an integer $r\geq1$, $t>0$, and $\beta=80$.  A partial coloring is a measurable map
$c:X\to\{0,1,\ldots,r\}$, where $0$ denotes an unassigned vertex.  Use the
shifts and polynomial choices of Section~\ref{sec:polynomials}, and let
$W:X\to(0,H_0]$ be a measurable weight, where $H_0<\infty$.  When a measurable
set of vertices has been deleted, take the principal restrictions of the
graph and matrix, retaining the original graph to measure distances.  At
absent vertices every rooted cost or transform is set equal to zero, and
the rooted measure is the zero measure.  On a
finite principal set $I$ we use the notation
\[
 P_I,\quad Q_{I,v}=P_{I\setminus\{v\}},\quad
 F_{I,v}=\frac{P_I}{Q_{I,v}},\quad G_{I,v}=\frac1{F_{I,v}},
 \quad \rho_{I,v},\quad m_{I,v},\quad d_{I,v}
\]
from the preceding section.  In particular, with roots counted with
multiplicity,
\[
 m_{I,v}=\int x_+\,\rho_{I,v}(dx),\qquad
 d_{I,v}=s(P_I)-s(Q_{I,v}),\qquad
 s(P)=\sum_{P(x)=0}x_+.
\]
We write $a_v=\tau_vW_v$ and set
\begin{equation}\label{inv:spectral-bound}
 L=(1+\beta t)H_0.
\end{equation}
All roots and rooted spectral measures in these finite calculations lie in
$[-L,L]$, including the virtual color choices at an unassigned vertex.

Measure invariance gives the mass transport identity
\begin{equation}\label{inv:mass-transport}
 \int_X\sum_{u\in[v]_{\mathcal R}} h(v,u)\,d\nu(v)
 =\int_X\sum_{v\in[u]_{\mathcal R}} h(v,u)\,d\nu(u)
\end{equation}
for nonnegative measurable $h$, and for integrable complex-valued $h$.
Indeed, partition $\mathcal R$ into countably many graphs of partial
one-to-one maps and change variables on each graph.

\subsection{Limits on bounded weighted graphs}

\begin{lemma}\label{lem:bounded-costs}
For the invariant measured graph data on the finite measure space
$(X,\nu)$ specified above, the rooted measures,
transforms, and costs have limits as finite principal subsets exhaust the
graph component of the root.  Denote them by
$\rho_v$, $G_v$, $F_v$, $m_v$, and $d_v$.  They are independent of the
exhaustion and are measurable in $v$.  At each present vertex, $\rho_v$ is a probability
measure supported on $[-L,L]$, and, for $z$ in the open upper half-plane,
\[
 G_v(z)=\int\frac{1}{z-x}\,\rho_v(dx),\qquad F_v(z)=\frac1{G_v(z)}.
\]
The same assertions hold after conditioning the color at one unassigned
vertex, with either its original shift or its increased shift.

For fixed $D_0,H_0,r,t$, the costs $m_v,d_v$ and the virtual shifted costs
$d_v^{\,i}$, $1\leq i\leq r$, are uniformly approximable by calculations on
a ball of finite radius about $v$.  The approximation radius is independent
of the particular graph, weights, partial coloring, and set of deleted
vertices.  These statements apply to a finite-measure restriction of an
invariant measured relation as well.
\end{lemma}

\begin{proof}
By Lemma~\ref{lem:rooted-locality}, each moment
$\int x^k\,\rho_{I,v}(dx)$ is determined by a bounded neighborhood of $v$,
and is unchanged once $I$ contains that neighborhood.  Compactness of the
probability measures on $[-L,L]$ and uniqueness from their moments therefore
give an exhaustion-independent weak limit $\rho_v$.  Its Cauchy transform
is the limit of $G_{I,v}$, locally uniformly on the upper half-plane.
Moreover,
\[
 -\operatorname{Im}G_v(z)
   =(\operatorname{Im}z)\int |z-x|^{-2}\,\rho_v(dx)>0,
\]
so this limit can be inverted.  Hence $F_{I,v}$ converges to $F_v$ there.
The same argument applies to the altered color law or shift at the root.
The finite quotients also have the form
\[
 F_{I,v}(z)-z-a_v=-\int\frac{1}{z-x}\,\kappa_{I,v}(dx),
\]
where the positive residue measure $\kappa_{I,v}$ is supported on
$[-L,L]$.  Its mass is at most $H_0^2$ by the second-moment formula in
Lemma~\ref{lem:rooted-locality}; the same bound holds for the conditional
color laws.  Compactness and uniqueness of Cauchy transforms therefore
give a limiting residue measure as well.  This also justifies taking
the local Laurent-coefficient limits of the quotients below.
Since $x_+$ is continuous on $[-L,L]$, the costs $m_{I,v}$ converge as
well.  Polynomial approximation on this interval, together with moment
locality, makes this convergence uniformly approximable by finite
neighborhood calculations.

Deletion increments require a different approximation, because their
definition involves the roots at every vertex of a finite set.  For a real
polynomial $f$, put
\[
 d_{I,v}[f]=\sum_{P_I(x)=0}f(x)
              -\sum_{Q_{I,v}(x)=0}f(x).
\]
The identity that the sum of the rooted measures counts all polynomial roots
gives
\begin{equation}\label{inv:polynomial-increment}
 d_{I,v}[f]
 =\sum_{u\in I}\left(
       \int f\,d\rho_{I,u}
       -1_{\{u\ne v\}}\int f\,d\rho_{I\setminus\{v\},u}
                    \right).
\end{equation}
Each summand is local.  It vanishes unless $u$ lies within a radius depending
only on $\deg f$ of $v$.  Thus the entire increment can be computed in a
sufficiently large finite neighborhood of $v$, independently of the ambient
finite set.  This defines its value on the infinite graph.

To pass from polynomials to $x_+$, for an absolutely continuous real
function $f$ on $[-L,L]$ use interlacing in the form
\begin{equation}\label{inv:ac-increment}
 d_{I,v}[f]
 =f(-L)+\int_{-L}^{L} f'(x)
           \bigl(n_{P_I}(x)-n_{Q_{I,v}}(x)\bigr)\,dx,
 \qquad
 0\leq n_{P_I}(x)-n_{Q_{I,v}}(x)\leq1.
\end{equation}
Here $n_P(x)$ counts roots strictly greater than $x$.  The formula follows
by writing $f(\lambda)=f(-L)+\int_{-L}^{\lambda}f'(x)\,dx$ for each root;
the difference in degrees is one.  Consequently
\begin{equation}\label{inv:increment-approximation}
 |d_{I,v}[f]-d_{I,v}[g]|
 \leq |f(-L)-g(-L)|+\|f'-g'\|_{L^1([-L,L])}.
\end{equation}
Approximate $f'$ in $L^1$ by polynomials and integrate them with prescribed
initial value $f(-L)$.  This gives polynomial approximants with both uniform
function error and the increment error in
Equation~\eqref{inv:increment-approximation} tending to zero.  Apply this
to $f(x)=x_+$.  The local polynomial increments are eventually constant,
and their approximation error is uniform in every finite set.  Hence
$d_{I,v}$ has the stated limit and uniform locality property.  The same
proof works for each of the finitely many virtual options at $v$.

Finally, finite balls in a measurable bounded-degree graph, their finite
matrix data, and their conjugacy-invariant polynomial calculations are
measurable.  The limiting constructions above are therefore measurable.
Deleting vertices or restricting the relation does not change the argument:
use the original graph for the locality radius and the retained vertices
for each principal calculation.
\end{proof}

\subsection{Integral change under ordered updates}

We next compare two global partial colorings, allowing deletion of vertices.
The order in which vertices are compared need not have a first element.
Only finitely many comparisons occur in any given polynomial local
calculation, which is enough for the following identity.

\begin{lemma}\label{lem:integral-change}
Under the bounded finite-measure hypotheses of
Lemma~\ref{lem:bounded-costs}, let the old and new data differ only by
changes of partial coloring or deletion at a measurable set $E_0$ of
vertices.  Fix a measurable real-valued function that is injective on
$E_0$ in each relation class, and order these vertices by its values.
At $u\in E_0$, use new data at all earlier vertices and old data at all
later vertices.  Let $\delta d_u$ be the new-minus-old difference of the
single-vertex deletion cost at $u$ in this environment, with cost zero
when $u$ is absent.  Then
\begin{equation}\label{inv:integral-change}
 \int_X(m_v^{\mathrm{new}}-m_v^{\mathrm{old}})\,d\nu(v)
       =\int_{E_0}\delta d_u\,d\nu(u).
\end{equation}
In particular, deletion alone cannot increase the integral of $m_v$.
\end{lemma}

\begin{proof}
An injective measurable point code supplies such an ordering if none has
already been specified.  Begin with a real polynomial $f$ in place of
$x_+$.  Write $m_v[f]=\int f\,d\rho_v$ at present vertices and zero at
absent ones.  For each root $v$, this quantity depends only on a fixed
finite-radius neighborhood.  List the finitely many vertices of $E_0$ in
that neighborhood in their specified order and telescope its change over
this list.  Each term agrees with the change produced at that vertex in
the global comparison environment: changes outside the neighborhood have
no effect on this polynomial rooted datum.  Thus no well-ordering or
infinite telescoping at an individual root is being assumed.

Integrate the resulting finite local sum and apply
Equation~\eqref{inv:mass-transport} to interchange the root and the changed
vertex.  For a fixed changed vertex $u$, the sum over affected roots equals
the change in $d_u[f]$.  Indeed, this can be checked on a sufficiently
large finite principal neighborhood using
Equation~\eqref{inv:polynomial-increment}: the old and new polynomials have
the same deletion polynomial at $u$.  If the vertex is deleted, the new
rooted datum there and the new deletion increment are both zero by
convention.  Bounded degree, bounded spectral support, and finite total
measure justify all integrals for this fixed polynomial.  We have proved
\[
 \int_X(m_v^{\mathrm{new}}[f]-m_v^{\mathrm{old}}[f])\,d\nu(v)
       =\int_{E_0}\delta d_u[f]\,d\nu(u).
\]

Choose the polynomial approximants to $x_+$ constructed in the proof of
Lemma~\ref{lem:bounded-costs}.  Uniform function approximation controls
the two rooted terms on the left; the $L^1$ derivative bound in
Equation~\eqref{inv:increment-approximation} controls both comparison
increments on the right.  Finite measure then permits passage to the
limit and proves Equation~\eqref{inv:integral-change}.  For a deletion
the comparison increment is $-d_u\leq0$, by interlacing.  This proves the
last assertion.
\end{proof}

\subsection{Deletion estimates independent of the largest weight}

The preceding arguments allow constants depending on $H_0$.  To pass later
to unbounded weights, we need a comparison on a fixed band that is uniform
as $H_0$ increases.  Cauchy transforms provide this comparison.  Besides
$G_v$, we must control the conditional color choices used in the
single-vertex assignment inequality.

At an unassigned vertex $v$, let $F_{v,i}^0$ be the quotient obtained by
conditioning its color to be $i$ while leaving its shift equal to $t$.
The permanent option has quotient
\[
 F_{v,i}(z)=F_{v,i}^0(z)+(\beta-1)tW_v.
\]
Define the auxiliary analytic field
\begin{equation}\label{inv:J-definition}
 J_i(v,z)=\bigl(F_{v,i}^0(z)-z-a_v\bigr)G_v(z),
       \qquad c(v)=0\text{ and }v\text{ present}.
\end{equation}
Set $J_i(v,z)=0$ at all other vertices. At present unassigned vertices the reciprocal transforms here are nonzero by
Lemma~\ref{lem:bounded-costs}.  The residue formula of
Lemma~\ref{lem:rooted-locality} shows that these analytic fields are
bounded on upper-half-plane compact sets when the root weight and its
neighbors' weight ratios are bounded.

\begin{lemma}\label{lem:rank-comparison}
Consider the invariant measured graph data on the finite measure space
$(X,\nu)$ specified above and assume, in addition, that neighboring
weights have ratios in $[K^{-1},K]$, where
$K\geq1$.  Fix $0<b_1<b_2<\infty$ and the band
\[
 \mathcal U=\{v:b_1\leq W_v\leq b_2\}.
\]
Delete a measurable set $E_{\mathrm{hi}}$ disjoint from $\mathcal U$, without
changing the coloring at retained vertices.  Let $\delta_{\mathrm{tot}}$
denote the difference between the original and the deleted configurations,
in either orientation.  For each compact set $\Omega$ in the open upper
half-plane there is a constant $C$ such that, for every $z\in\Omega$ and
$1\leq i\leq r$,
\begin{equation}\label{inv:rank-comparison}
 \int_{\mathcal U}|\delta_{\mathrm{tot}}G_v(z)|\,d\nu(v)
 +\int_{\mathcal U\cap\{c=0\}}
       |\delta_{\mathrm{tot}}J_i(v,z)|\,d\nu(v)
 \leq C\nu(E_{\mathrm{hi}}).
\end{equation}
Only band vertices present before deletion are included in these integrals.
The constant may depend on $\Omega,r,t,\beta,D_0,K,b_1,b_2$, but not on
$H_0$, the measure or its total mass, the partial coloring, the graph
itself, or the deleted set.
\end{lemma}

\begin{proof}
We may discard initially absent vertices from $E_{\mathrm{hi}}$: this
leaves both configurations unchanged and only decreases its measure.
Thus every deletion vertex below is present before the deletions begin.
We first prove bounds for one deletion in a finite principal set.  We then
pass these bounds to infinite bounded graphs and integrate them using an
analytic counterpart of the ordered-change identity.

\emph{A single deletion and the field $G$.}
Fix a finite principal set $I$, $u\in I$, and $z$ with
$y=\operatorname{Im}z>0$.  For arbitrary complex numbers
$\lambda_v$ satisfying $|\lambda_v|\leq1$, perturb the diagonal variables
to $z+\zeta\lambda_v$, $v\in I$.  On the disk $|\zeta|<y/2$ they remain
in the upper half-plane.  By Lemma~\ref{lem:polynomial-quotients}, the
quotient $P_I/P_{I\setminus\{u\}}$ in these variables has positive
imaginary part.  Its logarithm has a holomorphic branch with imaginary
part in $(0,\pi)$ throughout the disk.

We use the elementary derivative estimate that a holomorphic function on
a disk of radius $R$, with bounded imaginary part, has derivative at the
center bounded by a numerical constant times the oscillation of that
imaginary part divided by $R$.  To see this, restrict to a concentric
circle of radius $R/2$ and recover the first Taylor coefficient from the
first Fourier coefficient of the imaginary part.  Applied to the
logarithm just described, it bounds its derivative at zero by $C_0/y$.
Logarithmic differentiation of $P_I$ in a diagonal variable gives $G_{I,v}$,
so this derivative is
\[
 \sum_{v\in I}\lambda_v
       \bigl(G_{I,v}(z)-G_{I\setminus\{u\},v}(z)\bigr),
\]
where the second transform at $v=u$ is zero.  Choosing the phases of
$\lambda_v$ to make all summands nonnegative real yields
\begin{equation}\label{inv:single-G}
 \sum_{v\in I}
  |G_{I,v}(z)-G_{I\setminus\{u\},v}(z)|\leq C_0/y.
\end{equation}
This estimate contains no bound on the weights.

\emph{A single deletion and the conditional color field.}
Now assume $u\notin\mathcal U$, and take $\lambda_v=0$ except at
unassigned vertices in $\mathcal U$.  Perturb the off-diagonal parts of
both pair matrices of color $i$ by left and right diagonal scaling with
entries $1+\zeta\lambda_v$.  There is no conjugation of $\zeta$ in the
right scaling, and the fixed diagonal shift is left unchanged.  This
gives an analytic perturbation, though its matrices need not be Hermitian
for complex $\zeta$.

Let $D_\lambda=\diag(\lambda_v)$ and $D_W=\diag(\sqrt{W_v})$ on one
of the two relevant allowed-site sets.  With the sign suppressed, its
off-diagonal matrix is $C_i=D_W T D_W$, restricted to this set.  The
perturbation of this matrix equals
\[
 \zeta(D_\lambda C_i+C_iD_\lambda)
              +\zeta^2D_\lambda C_iD_\lambda.
\]
In the first row term only band rows and their graph neighbors occur.
The row weights are at most $b_2$, and these neighboring column weights
are at most $Kb_2$.  Factoring the diagonal weights on those row and
column sets and using $\|T\|\leq1$ gives
$\|D_\lambda C_i\|\leq\sqrt K b_2$.
The same argument applies to $C_iD_\lambda$.  In the quadratic term both
indices lie in the band, giving norm at most $b_2$.  The estimates apply
to every principal restriction.  Consequently the full matrix
perturbation has norm at most
\begin{equation}\label{inv:perturbation-bound}
             2|\zeta|\sqrt K b_2+|\zeta|^2b_2.
\end{equation}

Choose a disk of positive radius depending only on $y,b_2,K$ on which
Equation~\eqref{inv:perturbation-bound} is less than $y/2$.  The imaginary
parts of all perturbed pair matrices are then bounded above by
$(y/2)I$.  With diagonal variable still equal to $z$, the non-Hermitian
version of Lemma~\ref{lem:polynomial-quotients} shows that the quotient
$P_I/P_{I\setminus\{u\}}$ has positive imaginary part.  Once again its
logarithm has imaginary part in $(0,\pi)$, and the derivative estimate
is uniform in the upper weight bound $H_0$.

The derivative of $\log P_I$ in this perturbation is
\begin{equation}\label{inv:color-derivative}
               \frac2r\sum_{v\in I}\lambda_vJ_i(v,z).
\end{equation}
For completeness, consider the permutation expansion of an outcome
determinant.  At a target vertex $v$ whose selected color is $i$, a
nontrivial permutation cycle through $v$ uses one row factor and one
column factor, hence has derivative multiplier $2\lambda_v$.
A fixed point contributes $z+a_v$ and is unchanged because $T$ has zero
diagonal.  Conditional on color $i$ at $v$, the sum of the nontrivial
terms through $v$ is
\[
 P_{I\setminus\{v\}}(z)
       \bigl(F_{v,i}^0(z)-z-a_v\bigr).
\]
The probability of that color at an unassigned vertex is $1/r$.
Dividing by $P_I$ proves Equation~\eqref{inv:color-derivative}, including
the contribution of both signs.  The same formula holds after deleting
$u$.  Apply the derivative bound to the logarithm of the quotient and
choose the directions $\lambda_v$ as before.  We obtain
\begin{equation}\label{inv:single-J}
 \sum_{v\in I\cap\mathcal U\cap\{c=0\}}
       |J_i^{I}(v,z)-J_i^{I\setminus\{u\}}(v,z)|
       \leq C_1(y,b_2,K,r).
\end{equation}
The constants in Equations~\eqref{inv:single-G} and
\eqref{inv:single-J} can be chosen uniformly for $z$ in any fixed
upper-half-plane compact set.

\emph{Passage to bounded infinite graphs.}
Exhaust the graph component of $u$ by finite principal sets.  The
transforms converge by Lemma~\ref{lem:bounded-costs}, including the
conditional transforms in Equation~\eqref{inv:J-definition}.  Fatou's
Lemma therefore passes both single-deletion estimates to the infinite
graph.  Their sums can be taken over the entire relation class: there
is no effect on other graph components.  Moreover, the bounds hold in
every comparison environment formed by deleting other vertices first.

The single-deletion estimates are now uniform in the largest weight.
It remains to convert them into a comparison for a measurable set of
deletions.  Directly summing pointwise changes in an arbitrary infinite
order would not justify that step.  We instead establish the integrated
identity first for the local Laurent coefficients and then for the
analytic transforms.

\emph{The analytic integral change identity.}
Fix measurable complex directions $\lambda_v$, $|\lambda_v|\leq1$,
supported on the band targets: on $\mathcal U$ for $G$, and on
$\mathcal U\cap\{c=0\}$ for $J_i$.  Let $Z_v(z)$ denote the respective
field, extended by zero off those targets.  Order $E_{\mathrm{hi}}$
measurably and injectively along each class.  At a deletion vertex $u$,
let $\delta_u Z_v$ mean presence-minus-absence at $u$, with earlier
vertices already deleted and later vertices still present.  In the
following formula orient $\delta_{\mathrm{tot}}$ as
before-minus-after.  We claim that
\begin{equation}\label{inv:analytic-change}
 \int_X\lambda_v\delta_{\mathrm{tot}}Z_v(z)\,d\nu(v)
   =\int_{E_{\mathrm{hi}}}
       \left(\sum_{v\in[u]_{\mathcal R}}
                    \lambda_v\delta_u Z_v(z)\right)d\nu(u).
\end{equation}
The left side contains one rooted term, whereas the right side sums
over targets in the class of the deleted vertex.

To prove the identity, initially work near infinity in the upper
half-plane. Every target is present in both configurations, since the
deleted set is disjoint from the band and initially absent band vertices
are excluded. At such a target, the bounded spectral support gives
\[
 G_v(z)=z^{-1}\sum_{k\geq0}h_k(v)z^{-k},\qquad
                  |h_k(v)|\leq L^k.
\]
At a $J_i$ target, the corresponding expansion is
\[
 F_{v,i}^0(z)-z-a_v=-\sum_{k\geq0}b_k(v,i)z^{-k-1},\qquad
 |b_k(v,i)|\leq B_{v,i}L^k,
\]
where $B_{v,i}$ is the total mass of the limiting residue measure.
These bounds hold because the finite quotients are the negative Cauchy
transforms of their residue measures, all supported in $[-L,L]$.
The total residue mass is uniformly bounded in this bounded
setting by Lemma~\ref{lem:rooted-locality}.  These expansions persist in
the limit: each coefficient is eventually fixed by a finite
neighborhood, the coefficient bounds are uniform, and the series
converge to the limiting transforms for sufficiently large $|z|$.
The product defining $J_i$ therefore also has exponentially bounded
Laurent coefficients.  The coefficient at each order in either field
depends on a neighborhood of radius at most a constant times that
order, by Lemma~\ref{lem:rooted-locality} and multiplication of the
series.

For each individual coefficient the ordered local telescoping argument
from Lemma~\ref{lem:integral-change}, followed by
Equation~\eqref{inv:mass-transport}, proves the coefficient version of
Equation~\eqref{inv:analytic-change}.  The bounded degree ensures that
the number of affected roots at a fixed deletion vertex grows at most
exponentially in the coefficient order.  Combining this with the
coefficient bounds just obtained gives absolute summability of the
resulting series for sufficiently large $|z|$.  Finite measure then
allows the sums and integrals to be interchanged.  This proves
Equation~\eqref{inv:analytic-change} on a nonempty open region of the
upper half-plane.  The region may depend on $H_0$, which causes no
problem for the analytic continuation that follows.

Both sides of Equation~\eqref{inv:analytic-change} are holomorphic on
the full upper half-plane.  On the right, the sum over $v$ converges
absolutely at each $z$, by the single-deletion estimates.  Its analytic
partial sums, for example over increasing neighborhoods, are locally
bounded uniformly by those same estimates.  Normal-family convergence
therefore makes the sum holomorphic.  The outer integral is over the
finite measure set $E_{\mathrm{hi}}$, and the bounds are uniform on
compact sets, so it too is holomorphic.  On the left,
$|G_v(z)|\leq(\operatorname{Im}z)^{-1}$; the residue bound on the band
similarly bounds $J_i(v,z)$ on compact sets.  These estimates and the
finite band measure justify holomorphy of that integral.  The identity
theorem now proves Equation~\eqref{inv:analytic-change} everywhere in
the upper half-plane.

Finally fix $z$ and choose the measurable directions $\lambda_v$ so that
$\lambda_v\delta_{\mathrm{tot}}Z_v(z)=
|\delta_{\mathrm{tot}}Z_v(z)|$, with arbitrary value when the difference
is zero.  These are fixed directions for the just-proved identity; they
need not depend analytically on $z$.  Its right side is bounded by the
single-deletion constant times $\nu(E_{\mathrm{hi}})$.  Apply this
argument separately to $G$ and $J_i$ and add the estimates.  Their
constants were independent of $H_0$ and of the comparison environments,
which proves Equation~\eqref{inv:rank-comparison} with the stated
uniformity.
\end{proof}

The integral change identity controls the total positive cost under
color assignments and deletions.  The transform comparison gives a
different control: the effect of distant large weights on any fixed
weight band is bounded by the measure of the deleted vertices.  Together
these are the two estimates needed to define normalized costs on a
nonsingular base by passage through invariant restrictions of finite measure.
